% =============================================================================
%  InsaCloud — Les outils, brique par brique
%  Compilation :  tectonic -X compile technologies.tex
% =============================================================================
\documentclass[a4paper,11pt,oneside]{report}
% =============================================================================
%  Préambule du guide technique InsaCloud
% =============================================================================
\usepackage{fontspec}
% Noms de fichiers plutôt que noms système : la police vient du paquet
% tectonic, le document se recompile donc à l'identique sur n'importe quel poste.
\setmainfont{lmroman10-regular.otf}[
  BoldFont=lmroman10-bold.otf, ItalicFont=lmroman10-italic.otf,
  BoldItalicFont=lmroman10-bolditalic.otf]
\setsansfont{lmsans10-regular.otf}[
  BoldFont=lmsans10-bold.otf, ItalicFont=lmsans10-oblique.otf]
\setmonofont{lmmono10-regular.otf}[
  BoldFont=lmmonolt10-bold.otf, ItalicFont=lmmono10-italic.otf, Scale=0.92]
\usepackage[french]{babel}
\usepackage{graphicx,xcolor,booktabs,array,tabularx,longtable,multirow}
\usepackage{tikz}
\usetikzlibrary{positioning,arrows.meta,shapes.geometric,calc,fit,backgrounds,
                decorations.pathreplacing,patterns,shadows.blur}
\usepackage{fancyhdr,enumitem,microtype}
\usepackage[colorlinks=true,linkcolor=insared,urlcolor=insared,citecolor=insared]{hyperref}
\usepackage{caption}
\usepackage{listings}
\usepackage{tcolorbox}
\tcbuselibrary{skins,breakable}
\usepackage{titlesec}
\usepackage{emptypage}

% --- Couleurs ----------------------------------------------------------------
\definecolor{insared}{RGB}{200,16,46}
\definecolor{grisfond}{RGB}{247,248,250}
\definecolor{grisligne}{RGB}{216,220,226}
\definecolor{gristexte}{RGB}{91,97,107}
\definecolor{vertok}{RGB}{26,127,75}
\definecolor{vertfond}{RGB}{238,247,241}
\definecolor{orangepartiel}{RGB}{178,106,5}
\definecolor{orangefond}{RGB}{254,246,231}
\definecolor{rougeko}{RGB}{198,40,40}
\definecolor{bleufond}{RGB}{238,243,250}
\definecolor{bleutrait}{RGB}{43,91,158}

% --- Mise en page ------------------------------------------------------------
\captionsetup{font=small,labelfont={bf,color=insared}}
\setlist[itemize]{leftmargin=1.3em,itemsep=2pt,topsep=3pt,parsep=0pt}
\setlist[enumerate]{leftmargin=1.5em,itemsep=2pt,topsep=3pt,parsep=0pt}
\setlist[description]{leftmargin=1.3em,itemsep=3pt,topsep=3pt,parsep=0pt}

% LaTeX refuse par defaut qu'une figure occupe plus de 70 % d'une page et
% exige 20 % de texte : avec beaucoup de schemas, cela repousse les flottants
% et laisse des pages a moitie vides. On relache ces seuils.
\renewcommand{\topfraction}{0.92}
\renewcommand{\bottomfraction}{0.85}
\renewcommand{\textfraction}{0.06}
\renewcommand{\floatpagefraction}{0.75}
\setcounter{topnumber}{3}
\setcounter{bottomnumber}{2}
\setcounter{totalnumber}{4}

\pagestyle{fancy}
\fancyhf{}
\fancyhead[L]{\small\textcolor{gristexte}{\leftmark}}
\fancyhead[R]{\small\textcolor{gristexte}{InsaCloud}}
\fancyfoot[C]{\small\thepage}
\renewcommand{\headrulewidth}{0.4pt}
\renewcommand{\headrule}{\hbox to\headwidth{\color{grisligne}\leaders\hrule height \headrulewidth\hfill}}

\titleformat{\chapter}[display]
  {\normalfont\huge\bfseries}
  {\textcolor{insared}{\normalsize\MakeUppercase{\chaptertitlename} \thechapter}}
  {6pt}{\Huge}
\titlespacing*{\chapter}{0pt}{-20pt}{24pt}

% --- Extraits de code --------------------------------------------------------
\lstdefinestyle{insacloud}{
  basicstyle=\ttfamily\footnotesize,
  backgroundcolor=\color{grisfond},
  frame=single, rulecolor=\color{grisligne},
  breaklines=true, breakatwhitespace=false,
  postbreak=\mbox{\textcolor{gristexte}{$\hookrightarrow$}\space},
  showstringspaces=false,
  keywordstyle=\color{insared}\bfseries,
  commentstyle=\color{gristexte}\itshape,
  stringstyle=\color{vertok},
  numberstyle=\tiny\color{gristexte},
  xleftmargin=6pt, framexleftmargin=6pt, aboveskip=8pt, belowskip=8pt,
  extendedchars=true, inputencoding=utf8,
  literate=%
    {é}{{\'e}}1 {è}{{\`e}}1 {ê}{{\^e}}1 {ë}{{\"e}}1
    {à}{{\`a}}1 {â}{{\^a}}1 {ä}{{\"a}}1
    {î}{{\^i}}1 {ï}{{\"i}}1
    {ô}{{\^o}}1 {ö}{{\"o}}1
    {û}{{\^u}}1 {ù}{{\`u}}1 {ü}{{\"u}}1
    {ç}{{\c c}}1 {É}{{\'E}}1 {È}{{\`E}}1 {À}{{\`A}}1 {Ç}{{\c C}}1
    {œ}{{\oe}}1 {Œ}{{\OE}}1 {æ}{{\ae}}1
    {—}{{---}}1 {–}{{--}}1 {’}{{'}}1 {‘}{{'}}1
    {«}{{\guillemotleft}}1 {»}{{\guillemotright}}1
    {·}{{\textperiodcentered}}1 {→}{{$\rightarrow$}}1 {←}{{$\leftarrow$}}1
    {✔}{{\checkmark}}1 {°}{{\textdegree}}1 {€}{{\texteuro}}1
}
\lstset{style=insacloud}

\lstdefinelanguage{yamlplus}{
  keywords={true,false,null,name,hosts,become,roles,tasks,vars,when,loop,with_items,
            state,mode,owner,group,path,src,dest,notify,handlers,services,image,
            environment,volumes,networks,ports,depends_on,build,command,restart},
  keywordstyle=\color{insared}\bfseries,
  sensitive=false,
  comment=[l]{\#},
  morestring=[b]",
  morestring=[b]',
}

\lstdefinelanguage{dockerfile}{
  keywords={FROM,RUN,CMD,LABEL,EXPOSE,ENV,ADD,COPY,ENTRYPOINT,VOLUME,USER,
            WORKDIR,ARG,ONBUILD,STOPSIGNAL,HEALTHCHECK,SHELL,AS},
  keywordstyle=\color{insared}\bfseries,
  sensitive=true,
  comment=[l]{\#},
  morestring=[b]",
}

% --- Encadrés pédagogiques ---------------------------------------------------
\newtcolorbox{notion}[1]{
  breakable, enhanced, colback=bleufond, colframe=bleutrait,
  boxrule=0pt, leftrule=3pt, arc=2pt, left=8pt, right=8pt, top=6pt, bottom=6pt,
  fonttitle=\bfseries\small, title={Notion — #1},
  coltitle=bleutrait, titlerule=0pt, colbacktitle=bleufond,
  attach title to upper=\par\smallskip
}

\newtcolorbox{piege}[1]{
  breakable, enhanced, colback=orangefond, colframe=orangepartiel,
  boxrule=0pt, leftrule=3pt, arc=2pt, left=8pt, right=8pt, top=6pt, bottom=6pt,
  fonttitle=\bfseries\small, title={Piège — #1},
  coltitle=orangepartiel, titlerule=0pt, colbacktitle=orangefond,
  attach title to upper=\par\smallskip
}

\newtcolorbox{retenir}{
  breakable, enhanced, colback=vertfond, colframe=vertok,
  boxrule=0pt, leftrule=3pt, arc=2pt, left=8pt, right=8pt, top=6pt, bottom=6pt,
  fonttitle=\bfseries\small, title={À retenir},
  coltitle=vertok, titlerule=0pt, colbacktitle=vertfond,
  attach title to upper=\par\smallskip
}

% --- Raccourcis --------------------------------------------------------------
\newcommand{\ok}{\textcolor{vertok}{\textbf{Conforme}}}
\newcommand{\partiel}{\textcolor{orangepartiel}{\textbf{Partiel}}}
\newcommand{\plus}{\textcolor{insared}{\textbf{Au-delà}}}
\newcommand{\fichier}[1]{\textcolor{insared}{\texttt{#1}}}
\newcommand{\cmd}[1]{\texttt{#1}}

% Les paragraphes techniques enchaînent de longues suites en \texttt{} que TeX
% ne peut pas couper : sans marge de manœuvre, elles débordent dans la marge.
\emergencystretch=3em
\hbadness=6000
\tolerance=2000

% =============================================================================
%  Propre au document « technologies » : une fiche par outil
% =============================================================================
\usepackage[export]{adjustbox}

% Un bandeau d'ouverture pour chaque outil, avec son numéro de brique.
\newcommand{\outil}[2]{%
  \vspace{10pt}%
  \noindent\begin{tcolorbox}[enhanced, colback=grisfond, colframe=insared,
    boxrule=0pt, leftrule=4pt, arc=2pt, left=10pt, right=10pt,
    top=7pt, bottom=7pt, breakable]
  {\Large\bfseries #1}\hfill{\small\color{gristexte}#2}
  \end{tcolorbox}\vspace{2pt}}

% Les quatre questions auxquelles chaque fiche répond.
\newcommand{\quoi}{\par\vspace{5pt}\noindent{\color{insared}\textbf{C'est quoi}}\par\nopagebreak}
\newcommand{\pourquoi}{\par\vspace{7pt}\noindent{\color{insared}\textbf{Pourquoi celui-là}}\par\nopagebreak}
\newcommand{\role}{\par\vspace{7pt}\noindent{\color{insared}\textbf{Son rôle ici}}\par\nopagebreak}
\newcommand{\apport}{\par\vspace{7pt}\noindent{\color{insared}\textbf{Ce que ça a changé}}\par\nopagebreak}

% Capture recadrée : on ne montre que la zone dont parle le texte.
\newcommand{\gros}[3][0.8]{%
  \begin{figure}[htbp]\centering
  \includegraphics[width=#1\textwidth,frame]{figures/#2}
  \caption{#3}\end{figure}}

\renewcommand{\topfraction}{0.92}
\renewcommand{\bottomfraction}{0.85}
\renewcommand{\textfraction}{0.06}
\renewcommand{\floatpagefraction}{0.75}


\begin{document}

\begin{titlepage}
\centering
\vspace*{2.2cm}

{\Huge\bfseries InsaCloud\par}
\vspace{6pt}
{\color{insared}\rule{4cm}{2pt}\par}
\vspace{16pt}
{\LARGE Les outils, brique par brique\par}
\vspace{10pt}
{\large Chaque technologie : ce qu'elle est, pourquoi elle\par}
{\large et pas une autre, ce qu'elle fait ici\par}

\vspace{2cm}
\begin{minipage}{0.84\textwidth}
\color{gristexte}\small
Ce document ne raconte pas le projet : il prend les outils un par un.
Quarante-deux technologies, rangées en sept couches qui s'empilent. On part
d'un dossier vide et, à chaque brique posée, la plateforme sait faire une chose
de plus.

\medskip
Pour chacune, quatre questions, toujours les mêmes : \textbf{c'est quoi},
\textbf{pourquoi celle-là et pas une autre}, \textbf{son rôle exact ici}, et
\textbf{ce que son arrivée a changé}.
\end{minipage}

\vfill
{\color{gristexte}\small
Projet 3 --- \textit{Outils de déploiement de plateformes}\\
INSA Centre-Val de Loire, STI 4A FISE, 2026/2027 --- Dr.~Nadia Fettah\\[8pt]
\today\par}
\end{titlepage}

\pagenumbering{roman}
\setcounter{tocdepth}{1}
\tableofcontents
\cleardoublepage
\pagenumbering{arabic}

\chapter*{Comment lire ce document}
\addcontentsline{toc}{chapter}{Comment lire ce document}
\markboth{COMMENT LIRE CE DOCUMENT}{}

\section*{Le principe : un empilement}

On ne présente pas les outils par ordre alphabétique ni par importance. On les
pose dans l'ordre où il a fallu les poser, et chaque couche ne peut exister que
parce que la précédente est là.

\begin{figure}[htbp]
\centering
\resizebox{\textwidth}{!}{%
\begin{tikzpicture}[font=\small,
  c/.style={draw=grisligne,fill=grisfond,rounded corners=3pt,minimum width=11cm,
            minimum height=9mm,align=center},
  n/.style={font=\scriptsize\bfseries,insared}]
\foreach \y/\t/\o in {
  0/{Le socle}/{Git · Python · venv},
  1/{L'application web}/{Flask · Jinja2 · Werkzeug · Gunicorn · nginx · Tailwind},
  2/{Les données}/{SQLAlchemy · Alembic · SQLite · PostgreSQL · Fernet},
  3/{Les machines louées}/{Docker · supervisord · OpenSSH · ttyd · Xvnc · XFCE · noVNC},
  4/{Le pilotage du parc}/{subprocess · étiquettes · agent · Faucheur · systemd},
  5/{L'infrastructure}/{Vagrant · VMware · Ansible · Vault · inventaire dynamique},
  6/{La qualité}/{pytest · ruff · bandit · Trivy · Syft · GitHub Actions}} {
  \node[c] at (0,\y*1.15) {\textbf{\t}\\[1pt]\textcolor{gristexte}{\scriptsize\o}};
  \node[n] at (-6.3,\y*1.15) {\y};
}
\draw[-{Stealth[length=3mm]},insared,thick] (6.3,-0.45) -- (6.3,7.4)
  node[midway,right=4pt,insared,font=\scriptsize,align=left]
  {chaque couche\\suppose\\la précédente};
\end{tikzpicture}}
\caption{Les sept couches. On ne peut pas conteneuriser avant d'avoir une
application, ni automatiser le déploiement avant d'avoir quelque chose à
déployer.}
\end{figure}

\section*{Chaque outil, quatre questions}

Toutes les fiches ont la même forme. Vous pouvez lire le document en entier,
ou n'en consulter qu'une fiche.

\begin{description}
\item[C'est quoi] en une ou deux phrases, sans supposer qu'on connaît.
\item[Pourquoi celui-là] quelles étaient les autres possibilités, et ce qui a
  fait pencher. C'est la partie la plus utile : un outil se comprend par ce
  qu'on a écarté.
\item[Son rôle ici] précisément, dans ce projet, avec le fichier concerné.
\item[Ce que ça a changé] ce que la plateforme savait faire avant, et après.
\end{description}

\begin{retenir}
Si vous ne devez retenir qu'une chose par outil, retenez le \textbf{pourquoi
celui-là}. Savoir qu'on utilise Gunicorn ne vaut rien ; savoir pourquoi le
serveur de développement de Flask ne suffit pas, si.
\end{retenir}

\section*{Ce que la plateforme doit savoir faire}

Pour mesurer l'apport de chaque brique, il faut savoir où l'on va. La
plateforme finale doit :

\begin{enumerate}
\item accueillir des comptes utilisateurs et les authentifier ;
\item proposer un catalogue de systèmes Linux ;
\item créer une machine à la demande, en quinze secondes ;
\item en donner l'accès par SSH, par un terminal web et par un bureau
  graphique ;
\item la détruire toute seule à l'échéance ;
\item répartir les machines sur plusieurs serveurs ;
\item survivre à la panne de l'un d'eux ;
\item et se redéployer entièrement en une commande.
\end{enumerate}

Aucune de ces huit capacités n'est apportée par un seul outil. Chacune naît de
la rencontre de plusieurs briques --- et c'est précisément ce que ce document
cherche à rendre lisible.

\chapter{Couche 0 — Le socle}

Avant la première ligne de code métier, trois outils décident de ce que le
projet pourra devenir. On les oublie vite parce qu'ils ne « font » rien de
visible, et ce sont pourtant eux qu'on regrette quand ils manquent.

% -----------------------------------------------------------------------------
\outil{Git}{brique 1 · versionnement}

\quoi
Un système qui enregistre l'historique complet d'un dossier : qui a changé
quoi, quand, et pourquoi. Chaque enregistrement --- un \emph{commit} --- est
une photographie du projet entier, pas seulement des fichiers modifiés.

\pourquoi
L'alternative honnête, c'est de copier le dossier en
\cmd{projet-v2-final-vraiment}. Ça marche jusqu'au jour où l'on veut savoir
quand un bug est apparu, ou revenir trois jours en arrière sans perdre le
travail d'hier.

Il existe d'autres systèmes --- Subversion, Mercurial --- mais Git est celui que
toute l'industrie utilise, et surtout le seul que les plateformes
d'intégration continue savent lire nativement. Choisir autre chose aurait coupé
le projet de GitHub Actions, donc de toute la couche 6.

\role
Trente-cinq commits, chacun avec un message qui explique \emph{pourquoi} et pas
seulement \emph{quoi}. Les branches servent à isoler un travail en cours : la
branche \cmd{aziz-youssef-modif} porte l'intégration d'une contribution
extérieure sans toucher à \cmd{main} tant qu'elle n'est pas vérifiée.

Un fichier \fichier{.gitignore} tient hors du dépôt ce qui ne doit pas y
entrer : les clés (\cmd{.secret\_key}, \cmd{.vault\_key}), le fichier
\cmd{.env} des secrets de la pile, les environnements virtuels, et les archives
de travail des utilisateurs.

\apport
Sans Git, aucune des briques de la couche 6 n'existe : pas d'intégration
continue, pas de recherche de secrets dans l'historique, pas de déploiement
déclenché par un envoi de code. C'est la brique dont toutes les autres
dépendent indirectement.

% -----------------------------------------------------------------------------
\outil{Python 3.10}{brique 2 · langage}

\quoi
Le langage dans lequel est écrite toute la logique du projet : l'application
web, le démon de nettoyage, l'agent des nœuds, les scripts de génération de
schémas.

\pourquoi
Trois contraintes ont guidé le choix.

D'abord la \textbf{bibliothèque standard}. L'agent qui tourne sur chaque worker
ne doit rien avoir à installer : il utilise \cmd{urllib} pour les appels HTTP,
\cmd{json}, \cmd{ssl}, \cmd{subprocess}, et lit \cmd{/proc/meminfo}
directement. En Node.js ou en Go, il aurait fallu déployer un binaire ou un
environnement sur chaque nœud.

Ensuite l'\textbf{écosystème} : Flask, SQLAlchemy, Ansible lui-même sont en
Python. Rester dans un seul langage évite d'avoir deux chaînes d'outillage.

Enfin la \textbf{version} : 3.10 précisément, parce que la cible de
déploiement est Ubuntu 22.04, qui fournit cette version. Ce détail nous a
coûté une panne : nous avions épinglé une bibliothèque qui n'existait pas pour
3.10, les tests passaient sur nos postes et le déploiement échouait.

\begin{piege}{La version du poste n'est pas celle de la cible}
C'est l'erreur la plus banale et la plus coûteuse. La chaîne d'intégration
force désormais \cmd{PYTHON\_VERSION: 3.10} : elle teste avec le Python de la
machine de production, pas celui du développeur.
\end{piege}

\role
Environ 3\,500 lignes pour l'application, 1\,600 pour les tests, plus les
scripts d'outillage. Le code est séparé en modules dont chacun ignore les
couches au-dessus de lui : \fichier{docker\_ops.py} ne sait pas que Flask
existe, ce qui le rend testable sans serveur.

\apport
Un seul langage du navigateur au worker, et un agent déployable sans rien
installer.

% -----------------------------------------------------------------------------
\outil{venv et requirements épinglés}{brique 3 · reproductibilité}

\quoi
Un \emph{environnement virtuel} est un dossier contenant une copie isolée de
Python et de ses bibliothèques. Le fichier \fichier{requirements.txt} liste les
bibliothèques avec leur \textbf{version exacte} : \cmd{Flask==3.1.3}, pas
\cmd{Flask}.

\pourquoi
Sans isolation, installer une bibliothèque pour ce projet peut casser un autre
projet sur la même machine. Sans versions figées, le projet marche aujourd'hui
et casse demain parce qu'une dépendance a publié une version incompatible.

L'alternative moderne serait Poetry ou uv, plus confortables. On est resté sur
\cmd{venv} et \cmd{pip}, qui sont dans la bibliothèque standard : aucune
installation préalable n'est nécessaire sur la machine cible, et le rôle Ansible
peut créer l'environnement avec ce que fournit Ubuntu.

\role
Deux fichiers : \fichier{requirements.txt} pour ce qui tourne en production,
\fichier{requirements-dev.txt} pour l'outillage de tests et d'analyse. Le rôle
Ansible \cmd{webapp} installe exactement le premier sur le contrôleur --- pas
une liste recopiée à la main, ce qui a déjà été une source de divergence.

\apport
C'est ce qui permet à \cmd{pip-audit} d'exister (couche 6) : on ne peut
confronter des dépendances à une base de vulnérabilités que si l'on sait
exactement lesquelles on utilise.

% -----------------------------------------------------------------------------
\section*{Ce que la plateforme sait faire après la couche 0}

Rien d'utilisable. Mais le projet est désormais \emph{reproductible} : on peut
le cloner, créer l'environnement, et obtenir exactement les mêmes versions. Sans
cela, toutes les couches suivantes reposeraient sur du sable.

\chapter{Couche 1 — L'application web}

Six briques pour transformer du code Python en un site que des gens peuvent
utiliser. Chacune règle un problème que la précédente laissait ouvert.

\begin{figure}[htbp]
\centering
\resizebox{\textwidth}{!}{%
\begin{tikzpicture}[font=\scriptsize,
  b/.style={draw=grisligne,fill=grisfond,rounded corners=2pt,minimum height=13mm,
            minimum width=2.3cm,align=center},
  f/.style={-{Stealth[length=2mm]},gristexte}]
\node[b] (n) at (0,0) {Navigateur};
\node[b] (x) at (3.2,0) {\textbf{nginx}\\\textcolor{gristexte}{TLS, statiques}};
\node[b] (g) at (6.4,0) {\textbf{Gunicorn}\\\textcolor{gristexte}{4 ouvriers}};
\node[b] (fl) at (9.6,0) {\textbf{Flask}\\\textcolor{gristexte}{routes}};
\node[b] (j) at (12.8,0) {\textbf{Jinja2}\\\textcolor{gristexte}{gabarits}};
\draw[f] (n) -- node[above]{HTTPS} (x);
\draw[f] (x) -- node[above]{HTTP local} (g);
\draw[f] (g) -- node[above]{WSGI} (fl);
\draw[f] (fl) -- node[above]{rendu} (j);
\end{tikzpicture}}
\caption{Le chemin d'une requête. Chaque maillon existe parce que le suivant ne
sait pas faire son travail.}
\end{figure}

% -----------------------------------------------------------------------------
\outil{Flask}{brique 4 · le cadre web}

\quoi
Un \emph{micro}-framework : il fournit le routage (associer une URL à une
fonction), les gabarits et les sessions. Il ne fournit ni base de données, ni
administration, ni authentification --- on choisit soi-même.

\pourquoi
L'autre grand choix Python est \textbf{Django}, qui fournit tout : ORM,
interface d'administration, comptes utilisateurs, formulaires. Pour un projet
d'entreprise classique, Django va plus vite.

Ici, c'était un mauvais choix pour une raison précise : le sujet demande de
comprendre et de montrer chaque mécanisme. Avec Django, l'authentification est
une boîte noire qu'on active ; avec Flask, on écrit le hachage, la session, le
verrouillage après cinq échecs, et on peut les expliquer. Le projet est un
exercice d'ingénierie, pas une course à la livraison.

Deuxième raison : Django impose son ORM. Nous voulions pouvoir piloter la base
finement --- notamment l'index unique partiel qui garantit l'unicité des ports.

\role
\fichier{app.py} porte les routes HTML et les sessions ;
\fichier{api.py} porte les routes JSON. Une particularité : \textbf{les mêmes
URL servent les deux}, selon l'en-tête \cmd{Accept} de l'appelant. Il n'y a pas
un site d'un côté et une API de l'autre, mais une logique métier unique avec
deux rendus.

\begin{lstlisting}[language=Python,caption={\fichier{2\_webapp/app.py}}]
def veut_json() -> bool:
    # Un client qui accepte TOUT (Accept: */*, cas de curl sans en-tete)
    # n'est pas un client d'API : seul celui qui demande du JSON sans accepter
    # le HTML est traite comme tel.
    accepte = request.accept_mimetypes
    return accepte.accept_json and not accepte.accept_html
\end{lstlisting}

\apport
La plateforme devient atteignable depuis un navigateur. Et, grâce à la
négociation de contenu, elle est simultanément pilotable par \cmd{curl} sans
une ligne de code en plus --- ce qui satisfait l'exigence d'API du sujet sans
dédoubler le travail.

% -----------------------------------------------------------------------------
\outil{Jinja2}{brique 5 · les gabarits}

\quoi
Le moteur qui remplit des fichiers HTML avec des variables. Dans le gabarit on
écrit \verb|{{ machine.nom }}| et \verb|{% for m in machines %}| ; Jinja2
produit le HTML final.

\pourquoi
L'alternative serait de construire le HTML par concaténation de chaînes en
Python. Outre l'illisibilité, c'est une faille : une variable contenant
\cmd{<script>} serait injectée telle quelle dans la page. Jinja2
\textbf{échappe automatiquement} tout ce qu'il insère --- c'est une protection
contre les injections de script obtenue sans y penser.

Il vient avec Flask, donc aucun choix à faire ; mais le mentionner compte,
parce que l'échappement automatique est une des raisons pour lesquelles le
projet n'a aucune faille XSS.

\role
Les gabarits de \fichier{2\_webapp/templates/} : la page de connexion,
l'inscription, et le tableau de bord qui fait à lui seul 358 lignes. Les
conditions y servent à n'afficher une section que si elle a du contenu ---
« Archives de mon travail » n'apparaît que s'il y a des archives.

\apport
Des pages lisibles à écrire, et une classe entière d'attaques écartée par
défaut.

% -----------------------------------------------------------------------------
\outil{Werkzeug}{brique 6 · la plomberie HTTP}

\quoi
La bibliothèque sur laquelle Flask est construit. Elle gère le détail du
protocole HTTP, et fournit deux outils qu'on utilise directement :
le hachage des mots de passe et le \emph{middleware} \cmd{ProxyFix}.

\pourquoi
On ne l'a pas choisie : elle arrive avec Flask. Mais deux de ses fonctions sont
employées délibérément plutôt que réécrites.

\cmd{generate\_password\_hash} applique \textbf{pbkdf2-sha256} avec un sel
aléatoire et des milliers d'itérations. Écrire soi-même un hachage de mot de
passe est l'erreur classique : on choisit un algorithme trop rapide, on oublie
le sel, et la base devient exploitable en cas de fuite.

\cmd{ProxyFix} règle un problème qu'on ne voit qu'en production : derrière
nginx, toutes les requêtes semblent venir de \cmd{127.0.0.1}. Sans lui, notre
verrouillage anti-force-brute compterait toutes les tentatives comme venant de
la même adresse --- et bloquerait tout le monde au cinquième échec de n'importe
qui.

\role
Hachage à l'inscription et vérification à la connexion ; \cmd{ProxyFix} activé
dans \fichier{app.py} pour que l'adresse réelle du client remonte.

\apport
Des mots de passe stockés correctement, et un anti-force-brute qui vise la
bonne personne.

% -----------------------------------------------------------------------------
\outil{Gunicorn}{brique 7 · le serveur d'application}

\quoi
Un serveur WSGI : il reçoit les requêtes HTTP et les passe à l'application
Python. Il lance plusieurs processus ouvriers qui se partagent le travail.

\pourquoi
Flask embarque un serveur de développement, mais il traite \textbf{une requête
à la fois} et affiche lui-même un avertissement : ne pas utiliser en
production. Or créer un conteneur prend plusieurs secondes ; avec un seul
processus, un utilisateur qui loue une machine bloquerait tous les autres.

Les alternatives sont uWSGI (plus puissant, bien plus complexe à configurer) et
Waitress (plus simple, mais mono-processus par défaut). Gunicorn est le milieu
raisonnable : une ligne de commande suffit, et il redémarre les ouvriers qui
plantent.

\role
Quatre ouvriers, lancés par l'unité systemd du contrôleur :

\begin{lstlisting}[language=bash]
gunicorn --preload -w 4 -b 127.0.0.1:5000 app:app
\end{lstlisting}

\begin{piege}{Quatre processus qui démarrent en même temps}
Chacun importe l'application, qui crée les tables si elles manquent. Les quatre
le font simultanément, et parfois l'un échoue avec \cmd{table users already
exists}. Sur un poste mono-processus, le bug n'apparaît jamais.

Deux correctifs : la création tolère désormais l'erreur « existe déjà », et
l'option \cmd{--preload} fait importer l'application \emph{une fois} avant de
dupliquer le processus. Nous avons reproduit la panne avec huit processus
simultanés --- six sur huit démarraient ; après correction, huit sur huit.
\end{piege}

\apport
La plateforme supporte plusieurs utilisateurs simultanés. Avant Gunicorn, une
location en cours gelait le site pour tout le monde.

% -----------------------------------------------------------------------------
\outil{nginx}{brique 8 · le proxy inverse}

\quoi
Un serveur web placé \emph{devant} l'application. Le navigateur lui parle ; il
relaie vers Gunicorn.

\pourquoi
Gunicorn sait répondre à du HTTP, mais ne sait pas faire trois choses
essentielles : terminer le TLS, servir efficacement les fichiers statiques, et
absorber les requêtes malformées.

L'alternative est Apache, qui sait tout faire aussi mais consomme davantage et
se configure plus lourdement. Caddy aurait apporté les certificats
automatiques, ce qui n'a pas de sens sur un réseau privé sans nom de domaine
public.

\role
Le rôle Ansible \cmd{reverse\_proxy} installe nginx avec :

\begin{itemize}
\item \textbf{TLS 1.3} et le certificat du nœud ;
\item les fichiers statiques servis directement, sans réveiller Python ;
\item les en-têtes de sécurité : \cmd{Strict-Transport-Security},
  \cmd{X-Content-Type-Options}, \cmd{X-Frame-Options} ;
\item \cmd{X-Forwarded-For}, que \cmd{ProxyFix} relira côté Flask.
\end{itemize}

\apport
Le site passe en HTTPS, les images et le CSS ne coûtent plus de temps Python, et
l'application n'est plus exposée directement au réseau.

% -----------------------------------------------------------------------------
\outil{Tailwind CSS, précompilé}{brique 9 · l'apparence}

\quoi
Une bibliothèque de style où l'on compose l'apparence avec de petites classes
dans le HTML, plutôt qu'en écrivant du CSS à part.

\pourquoi
L'usage courant est de charger Tailwind depuis un CDN : une ligne, et ça marche.
C'était impossible ici.

La politique de sécurité de contenu du projet est \cmd{default-src 'self'} :
\textbf{aucune ressource externe} n'a le droit d'être chargée. Un CDN est
précisément une ressource externe. Il a donc fallu \textbf{précompiler} Tailwind
en un fichier CSS servi par nginx, et héberger les polices localement.

\role
\fichier{2\_webapp/static/css/tailwind.css} est le résultat de la compilation ;
\fichier{tailwind.src.css} en est la source. Les polices IBM Plex Sans sont dans
\fichier{static/fonts/}.

\apport
Une interface soignée \emph{sans} affaiblir la politique de sécurité. C'est le
genre de compromis où le choix facile (le CDN) aurait coûté une faille entière :
avec un CDN autorisé, une injection de script redevient exploitable.

\gros[0.55]{zoom-connexion.png}{La page de connexion. Tout le style vient d'un
fichier CSS local : la politique de sécurité interdit toute ressource externe.}

% -----------------------------------------------------------------------------
\section*{Ce que la plateforme sait faire après la couche 1}

Elle a un site en HTTPS, capable de servir plusieurs utilisateurs, avec des
comptes et des sessions. Mais elle ne retient rien : à chaque redémarrage, tout
disparaît. C'est l'objet de la couche suivante.

\chapter{Couche 2 — Les données}

L'application a des comptes et des sessions, mais elle oublie tout au
redémarrage. Cinq briques pour qu'elle se souvienne --- et pour que ce dont elle
se souvient reste secret.

% -----------------------------------------------------------------------------
\outil{SQLAlchemy 2.0}{brique 10 · l'ORM}

\quoi
Un \emph{Object-Relational Mapper} : il fait correspondre une classe Python à
une table SQL. On écrit des objets, il écrit le SQL.

\pourquoi
L'alternative est d'écrire le SQL à la main. Trois raisons l'ont écartée.

\textbf{Les injections SQL.} Écrire
\cmd{f"SELECT * FROM users WHERE nom='\{saisie\}'"} laisse un utilisateur nommé
\cmd{' OR '1'='1} lire toute la table. L'ORM passe toujours par des paramètres
liés : l'injection devient structurellement impossible, pas seulement évitée par
discipline.

\textbf{Deux moteurs.} Le projet tourne sur SQLite en développement et sur
PostgreSQL en production. Leur SQL diffère sur les dates, les types, les
clauses. L'ORM écrit le dialecte qui convient.

\textbf{La vérification.} Une faute de frappe dans un nom de colonne devient une
erreur Python immédiate, au lieu d'une erreur SQL découverte en production.

On aurait pu prendre l'ORM de Django, mais il impose Django. Celui de SQLAlchemy
est le standard du monde Python hors Django, et sa version 2.0 apporte une
écriture typée que les analyseurs comprennent.

\role
\fichier{2\_webapp/models.py} décrit les sept tables. Chaque colonne est annotée
pour que \cmd{ruff} et l'éditeur sachent son type :

\begin{lstlisting}[language=Python]
class Rental(Base):
    __tablename__ = "rentals"
    id: Mapped[int] = mapped_column(Integer, primary_key=True)
    user_id: Mapped[int] = mapped_column(ForeignKey("users.id"), nullable=False)
    end_time: Mapped[datetime] = mapped_column(DateTime, nullable=False)
    __table_args__ = (Index("idx_rentals_end", "status", "end_time"),)
\end{lstlisting}

L'index de la dernière ligne n'est pas décoratif : c'est exactement la requête
que le Faucheur exécute toutes les dix secondes.

\apport
La plateforme se souvient. Et une famille entière d'attaques --- l'injection
SQL --- disparaît du périmètre.

% -----------------------------------------------------------------------------
\outil{Alembic}{brique 11 · les migrations}

\quoi
L'outil qui fait évoluer le schéma d'une base \emph{sans perdre les données}.
Chaque changement devient un fichier numéroté qui connaît son prédécesseur ; la
base retient où elle en est.

\pourquoi
Tant qu'on développe, on efface la base et on recommence. Dès qu'il y a des
données réelles --- des comptes, des locations en cours --- c'est impossible.

Sans Alembic, la seule option serait d'écrire les \cmd{ALTER TABLE} à la main et
de se souvenir lesquels ont été appliqués où. C'est exactement le genre de
mémoire humaine qui échoue au pire moment.

\role
\fichier{2\_webapp/migrations/} contient deux révisions : le schéma initial, puis
l'ajout de la table \cmd{archives}. La chaîne d'intégration exécute
\cmd{alembic check}, qui échoue si les modèles et les migrations divergent ---
impossible d'oublier de générer une migration.

\begin{piege}{SQLite ne sait pas modifier une colonne}
PostgreSQL accepte \cmd{ALTER TABLE ... ALTER COLUMN}. SQLite, non : il faut
créer une table neuve, recopier, supprimer, renommer. Alembic sait le faire,
mais seulement si on le lui demande par \cmd{render\_as\_batch=True}. Sans cette
option, les migrations passent sur un moteur et échouent sur l'autre.
\end{piege}

\apport
Le schéma peut évoluer en production. C'est ce qui a permis d'ajouter la
sauvegarde du travail --- une table de plus --- sur une base contenant déjà des
comptes et un historique.

% -----------------------------------------------------------------------------
\outil{SQLite}{brique 12 · la base de développement}

\quoi
Une base de données contenue dans un seul fichier, sans serveur à lancer.

\pourquoi
Pour développer et pour les tests, démarrer un serveur PostgreSQL est une
friction : il faut l'installer, le configurer, le lancer. SQLite est déjà dans
Python.

Le projet l'utilise en mode \textbf{WAL} (\emph{write-ahead logging}), qui
permet à plusieurs processus de lire pendant qu'un autre écrit --- nécessaire
puisque l'application et le Faucheur touchent la même base.

\role
Base par défaut quand \cmd{DATABASE\_URL} n'est pas défini. Les 151 tests
l'utilisent : chacun part d'un fichier vierge dans un dossier temporaire, ce qui
garantit qu'aucun test n'hérite de l'état d'un autre.

\apport
Les tests tournent en vingt secondes sur n'importe quelle machine, sans service
à installer. C'est ce qui rend l'intégration continue légère.

% -----------------------------------------------------------------------------
\outil{PostgreSQL 16}{brique 13 · la base de production}

\quoi
Un serveur de base de données complet, conçu pour plusieurs clients simultanés.

\pourquoi
SQLite suffit à un processus, pas à quatre ouvriers Gunicorn plus le Faucheur
qui écrivent en même temps. Son verrouillage porte sur le fichier entier : sous
charge, les écritures s'attendent.

MySQL aurait fait l'affaire ; PostgreSQL a été préféré pour deux raisons
concrètes ici : les \textbf{index partiels} --- la clause \cmd{WHERE} sur un
index unique, qui garantit l'unicité des ports sans bloquer leur
réutilisation --- et le respect strict des types.

\role
Service \cmd{db} de la pile compose, image \cmd{postgres:16-alpine}. Son port
\textbf{n'est pas publié} : seuls les services du réseau interne y accèdent. Il
démarre avec \cmd{--data-checksums}, qui détecte une corruption silencieuse du
disque.

\apport
La plateforme supporte la concurrence réelle. Et l'index unique partiel devient
une garantie du moteur, pas une intention du code.

% -----------------------------------------------------------------------------
\outil{cryptography · Fernet}{brique 14 · le coffre}

\quoi
Fernet est une recette de chiffrement symétrique authentifié : elle combine AES
en mode CBC et un HMAC. On ne peut ni lire sans la clé, ni modifier sans que
cela se voie.

\pourquoi
Le problème est particulier : le mot de passe root d'une machine louée doit
pouvoir être \textbf{réaffiché}. On ne peut donc pas le hacher --- le hachage est
à sens unique. Il faut un chiffrement réversible.

Écrire soi-même du chiffrement est l'erreur classique : on choisit mal le mode,
on réutilise un vecteur d'initialisation, on oublie l'authentification. Fernet
n'a \textbf{aucun paramètre à choisir} : c'est précisément son intérêt pour qui
n'est pas cryptographe.

\begin{notion}{Hacher ou chiffrer : deux besoins, deux outils}
Le mot de passe d'un \textbf{compte} est \textbf{haché} : on ne doit jamais
pouvoir le retrouver, seulement vérifier qu'un essai donne la même empreinte.

Le mot de passe \textbf{root d'une machine louée} est \textbf{chiffré} : il doit
pouvoir être réaffiché à son propriétaire.

Hacher le second le rendrait inutilisable ; chiffrer le premier serait une
faute.
\end{notion}

\role
\fichier{2\_webapp/crypto.py} charge deux clés \textbf{distinctes} : celle qui
signe les sessions Flask, et celle qui chiffre les mots de passe root. Les
séparer garantit qu'un vol de la première ne donne pas accès aux machines.

Les clés viennent de l'environnement --- Ansible les injecte depuis son Vault ---
ou, à défaut, d'un fichier local en permissions \cmd{0600}.

\gros[0.86]{zoom-secret.png}{Le mot de passe root, affiché une seule fois. Après
cet instant il n'existe plus qu'en base, chiffré.}

\gros[0.86]{zoom-coffre.png}{Le coffre rouvert après re-saisie du mot de passe
de compte : les accès redeviennent visibles cinq minutes, puis se remasquent.}

\begin{piege}{Une session Flask est signée, pas chiffrée}
Notre première version déposait le mot de passe root dans le cookie de session,
le temps de l'afficher. Or le contenu d'un cookie de session Flask est encodé en
base64 : \textbf{n'importe qui peut le lire} sans la clé. La signature empêche
seulement de le modifier.

On ne stocke plus que l'identifiant de la machine ; le secret est relu en base
et déchiffré au moment de l'affichage, après vérification du propriétaire.
\end{piege}

\apport
Un secret qu'on peut revoir sans jamais le stocker en clair, et qui ne transite
plus par le navigateur.

% -----------------------------------------------------------------------------
\section*{Ce que la plateforme sait faire après la couche 2}

Elle retient les comptes, les locations et les machines ; elle sait faire
évoluer son schéma sans perdre de données ; et elle garde des secrets qu'elle
peut réafficher. Il lui manque l'essentiel : les machines à louer.

\chapter{Couche 3 — Les machines louées}

C'est le cœur du produit. Sept briques pour fabriquer ce que l'utilisateur
achète : une machine Linux à lui, joignable de trois façons.

% -----------------------------------------------------------------------------
\outil{Docker}{brique 15 · la conteneurisation}

\quoi
Un moteur qui exécute des processus isolés les uns des autres, en réutilisant le
noyau de la machine hôte. Un \emph{conteneur} se croit seul au monde, mais ne
porte pas de système d'exploitation complet.

\pourquoi
La vraie alternative était la machine virtuelle --- VirtualBox, KVM. Une VM
offre une isolation plus forte, puisqu'elle a son propre noyau. Deux chiffres
ont tranché :

\begin{center}
\small
\begin{tabularx}{0.9\textwidth}{@{}lXX@{}}
\toprule
& \textbf{Machine virtuelle} & \textbf{Conteneur} \\
\midrule
Démarrage & 30 à 60 secondes & moins d'une seconde \\
Mémoire minimale & plusieurs centaines de Mo & celle des processus \\
Densité sur un serveur & quelques dizaines & plusieurs centaines \\
\bottomrule
\end{tabularx}
\end{center}

Pour une location \textbf{à la minute}, faire attendre l'utilisateur une minute
est absurde. Le sujet tranche d'ailleurs explicitement : « des conteneurs dans
notre cas ».

\begin{notion}{Ce qu'un conteneur est réellement}
Trois mécanismes du noyau Linux, pas un objet magique :

\textbf{Les espaces de noms} cloisonnent ce que le processus \emph{voit} : son
propre arbre de processus, ses propres interfaces réseau, sa propre vision du
disque. Le port 22 du conteneur n'est pas le port 22 de l'hôte.

\textbf{Les groupes de contrôle} limitent ce qu'il \emph{consomme} : mémoire,
processeur, nombre de processus.

\textbf{Les capacités} découpent les pouvoirs de root en une quarantaine de
droits indépendants, qu'on accorde un par un.
\end{notion}

\role
Six images, et une commande imposée par le sujet pour créer chaque machine :

\begin{lstlisting}[language=bash]
docker run -d --restart=always -p <port>:22 <image>
\end{lstlisting}

Le projet y ajoute le durcissement : \cmd{--cap-drop ALL} puis les dix capacités
strictement nécessaires, \cmd{no-new-privileges}, une limite de processus, de
CPU, et la mémoire en réservation \emph{et} en limite.

\apport
La plateforme sait livrer une machine en quinze secondes. C'est la brique qui
rend le produit possible.

% -----------------------------------------------------------------------------
\outil{Dockerfile}{brique 16 · la recette des images}

\quoi
Un fichier texte qui décrit, instruction par instruction, comment construire
une image. Chaque instruction crée une \emph{couche} empilée sur la précédente.

\pourquoi
On pourrait créer une image en lançant un conteneur, en l'installant à la main,
puis en figeant le résultat. Personne ne saurait plus ce qu'il contient, ni
comment le reproduire.

Le Dockerfile est l'infrastructure comme code appliquée aux images : versionné,
relu, et vérifiable par un analyseur.

\role
Quatre fichiers pour six images. Le mode bureau n'est pas un fichier séparé :
c'est un paramètre de construction.

\begin{lstlisting}[language=dockerfile]
FROM ubuntu:22.04
ARG DESKTOP=0
RUN apt-get update \
    && apt-get install -y --no-install-recommends openssh-server supervisor \
    && rm -rf /var/lib/apt/lists/*
COPY entrypoint.sh /usr/local/bin/entrypoint.sh
CMD ["/usr/local/bin/entrypoint.sh"]
\end{lstlisting}

\begin{piege}{Nettoyer dans une autre couche ne sert à rien}
Si l'on installe des paquets dans un \cmd{RUN} et qu'on efface le cache d'apt
dans le \cmd{RUN} suivant, le cache reste dans l'image : une couche ne peut pas
alléger celle du dessous. Il faut enchaîner les deux dans la \emph{même}
instruction. Sur nos images, cela représente une cinquantaine de mégaoctets.
\end{piege}

\begin{piege}{Un mot de passe par défaut dans l'image}
Notre première version contenait \cmd{ENV ROOT\_PASSWORD=insacloud}. L'analyseur
\cmd{hadolint} l'a signalé, et il avait raison : ce mot de passe aurait été le
même sur \textbf{toutes} les machines louées, et visible par quiconque fait
\cmd{docker history}. L'entrypoint en tire désormais un au hasard si la
plateforme n'en fournit pas.
\end{piege}

\apport
Six images reproductibles, construites par la chaîne d'intégration à chaque
envoi, et analysées avant d'être publiées.

% -----------------------------------------------------------------------------
\outil{supervisord}{brique 17 · plusieurs services dans un conteneur}

\quoi
Un gestionnaire de processus : il en démarre plusieurs, les surveille, et
relance celui qui meurt.

\pourquoi
C'est la brique qui résout une contradiction. Un conteneur exécute \emph{un}
processus, et s'arrête quand celui-ci se termine. Or une machine louée doit faire
tourner cinq services : SSH, le terminal web, le serveur d'affichage, le bureau,
et la passerelle VNC.

Les alternatives : lancer un script shell qui démarre tout en arrière-plan ---
mais alors rien ne surveille, et un service mort le reste. Ou \cmd{systemd} dans
le conteneur --- lourd, et il réclame des privilèges qu'on vient justement de
retirer.

\begin{notion}{Pourquoi le PID 1 n'est pas un processus comme un autre}
Le premier processus d'un conteneur hérite d'une obligation : adopter les
processus orphelins et récupérer leur code de sortie. Sans cela, les processus
terminés s'accumulent à l'état \emph{zombie} et finissent par saturer la table.

Un \cmd{bash} lancé comme PID 1 ne fait pas ce travail. \cmd{supervisord}, si.
\end{notion}

\role
L'entrypoint écrit une configuration adaptée au mode choisi, puis se
\textbf{remplace} par supervisord :

\begin{lstlisting}[language=bash]
exec "$(command -v supervisord)" -c "$CONF"
\end{lstlisting}

Le mot \cmd{exec} compte : sans lui, le shell resterait PID 1 et ne relaierait
pas les signaux. Un \cmd{docker stop} deviendrait un arrêt brutal au bout du
délai de grâce.

\apport
C'est le premier des trois niveaux de résilience du projet : si un service meurt
\emph{dans} la machine, il est relancé sans que l'utilisateur s'en aperçoive.

% -----------------------------------------------------------------------------
\outil{OpenSSH}{brique 18 · l'accès demandé par le sujet}

\quoi
Le serveur SSH : il donne un terminal distant chiffré, authentifié par mot de
passe ou par clé.

\pourquoi
Il n'y avait pas de choix à faire : c'est l'accès que le sujet impose, et c'est
le standard universel. La question intéressante est plutôt \emph{comment} on
l'expose.

Le conteneur écoute sur son propre port 22, et Docker le publie sur un port de
l'hôte dans la plage 8000--9000. C'est ce qui permet de faire tourner des
dizaines de serveurs SSH sur une seule machine : chacun croit avoir le port 22.

\role
Deux modes d'authentification. Par défaut, le mot de passe root généré. Mais si
l'utilisateur a enregistré sa clé publique sur son compte, la machine
\textbf{n'accepte plus que} cette clé --- l'authentification par mot de passe
est coupée.

\begin{piege}{Le port est publié avant que le service écoute}
Docker publie le port dès la création du conteneur, mais \cmd{sshd} met une
seconde ou deux à démarrer. Un test qui se connecte immédiatement échoue de
façon intermittente.

La correction n'est pas d'attendre « assez longtemps » --- une durée fixe est
toujours trop courte un jour --- mais de réessayer jusqu'à ce que le service
réponde, avec une limite.
\end{piege}

\apport
L'accès exigé par le cahier des charges, avec une option plus sûre que le mot
de passe pour qui le souhaite.

% -----------------------------------------------------------------------------
\outil{ttyd}{brique 19 · le terminal dans le navigateur}

\quoi
Un petit serveur qui expose un terminal dans une page web, par WebSocket.

\pourquoi
Le sujet ne le demande pas. Il répond à un problème réel : un utilisateur sur un
poste de salle de TP n'a pas toujours de client SSH, ni le droit d'en installer.

Les alternatives --- Shellinabox, abandonné ; Wetty, en Node.js donc une chaîne
d'outillage de plus --- ont été écartées. ttyd est un binaire unique, en C, de
quelques centaines de kilooctets.

\role
Lancé par supervisord sur le port 7681 du conteneur, publié sur l'hôte. Détail
qui compte : il ne lance pas un shell, il lance \cmd{login}.

\begin{lstlisting}[language=bash]
ttyd -p 7681 -W /bin/sh -c "cat /etc/issue; exec $(command -v login)"
\end{lstlisting}

Le port ne donne donc pas un accès ouvert : il présente une invite qui réclame
l'identifiant et le mot de passe.

\gros[0.8]{zoom-motd.png}{Ce que montre le port du terminal web : une invite de
connexion, pas un shell.}

\gros[0.8]{zoom-bienvenue.png}{Après authentification, le bandeau rappelle que
la machine est louée pour une durée limitée.}

\apport
Un accès sans rien installer. C'est le premier des ajouts au-delà du sujet.

% -----------------------------------------------------------------------------
\outil{Xvnc, XFCE et noVNC}{brique 20 · le bureau graphique}

\quoi
Trois pièces qui s'emboîtent. \textbf{Xvnc} est un serveur d'affichage qui
dessine dans une mémoire plutôt que sur un écran. \textbf{XFCE} est
l'environnement de bureau qui s'y affiche. \textbf{noVNC} est une passerelle qui
traduit le protocole VNC en WebSocket, pour qu'un navigateur puisse l'afficher.

\pourquoi
XFCE plutôt que GNOME ou KDE : il tourne correctement dans un gigaoctet de
mémoire, là où les autres en réclament deux ou trois. Pour une machine louée,
c'est la différence entre trois utilisateurs par serveur et dix.

noVNC plutôt qu'un client VNC à installer : même raisonnement que pour ttyd ---
l'utilisateur ne doit rien installer.

\role
Trois services de plus sous supervisord, en mode bureau seulement. Le serveur
d'affichage \textbf{n'écoute que sur la boucle locale} : on ne peut l'atteindre
qu'à travers la passerelle, qui est chiffrée quand un certificat est monté.

Le mot de passe VNC est limité à huit caractères par le protocole : ce sont les
huit premiers du mot de passe root, ce que l'interface explique.

\gros[0.86]{zoom-bureau.png}{Le bureau XFCE d'une machine Ubuntu louée, affiché
dans le navigateur. Cinq processus cohabitent dans ce conteneur.}

\apport
Le second ajout au-delà du sujet, et le plus spectaculaire en démonstration : un
bureau Linux complet, dans un onglet, sans rien installer.

% -----------------------------------------------------------------------------
\outil{Le script d'entrée}{brique 21 · ce qui assemble la machine}

\quoi
\fichier{1\_docker/entrypoint.sh}, 166 lignes de shell. C'est le programme que
Docker lance quand la machine démarre.

\pourquoi
Il fallait un endroit unique qui adapte la machine à ce que l'utilisateur a
demandé. Mettre cette logique dans le Dockerfile aurait obligé à six images
indépendantes ; la mettre dans l'application aurait dispersé le savoir.

\role
Trois tâches, dans l'ordre :

\begin{enumerate}
\item poser le mot de passe root et, s'il y en a une, la clé SSH --- auquel cas
  il \textbf{coupe} l'authentification par mot de passe ;
\item écrire la configuration supervisord correspondant au mode ;
\item se remplacer par supervisord.
\end{enumerate}

Il sert aussi d'outil : la plateforme l'appelle pour changer le mot de passe
d'une machine à chaud, sans la redémarrer.

\begin{lstlisting}[language=bash]
docker exec <conteneur> /usr/local/bin/entrypoint.sh setpass <nouveau>
\end{lstlisting}

\apport
Une seule recette pour six variantes, et la rotation du mot de passe sans
interruption de service.

% -----------------------------------------------------------------------------
\section*{Ce que la plateforme sait faire après la couche 3}

Elle fabrique de vraies machines Linux, accessibles de trois façons, durcies, et
qui se réparent toutes seules en interne. Mais quelqu'un doit encore décider
\emph{quand} les créer et \emph{quand} les détruire --- et sur quel serveur.

\chapter{Couche 4 — Le pilotage du parc}

Les machines existent, mais rien ne décide quand les créer, où les placer, ni
quand les détruire. Sept briques pour transformer un démon Docker en un service
de location.

% -----------------------------------------------------------------------------
\outil{subprocess}{brique 22 · piloter Docker depuis Python}

\quoi
Le module de la bibliothèque standard qui lance un programme externe et récupère
sa sortie. Ici, il lance la commande \cmd{docker}.

\pourquoi
Il existe une bibliothèque Python officielle, \cmd{docker-py}, qui parle à
l'API Docker sans passer par la ligne de commande. Elle est plus élégante.

On l'a écartée pour une raison de \textbf{conformité et de lisibilité}. Le sujet
impose une commande précise :

\begin{center}
\cmd{docker run -d --restart=always -p <port>:22 <image>}
\end{center}

Avec un SDK, cette commande n'apparaît nulle part : elle devient un appel de
méthode avec des paramètres nommés. Avec \cmd{subprocess}, chaque option se lit
telle quelle dans le code, et l'on peut montrer à la prof la ligne exacte qui
part.

\role
Une seule fonction centralise tous les appels, dans
\fichier{2\_webapp/docker\_ops.py} :

\begin{lstlisting}[language=Python]
def run_docker(*args, worker="local", timeout=DOCKER_TIMEOUT) -> str:
    cmd = [*wk.docker_command(worker), *args]
    # Appel sans shell : les arguments sont passes en liste, donc jamais
    # re-interpretes par un interpreteur de commandes.
    res = subprocess.run(cmd, capture_output=True, text=True, timeout=timeout)
\end{lstlisting}

\begin{piege}{\texttt{shell=True} est une porte ouverte}
Écrire \cmd{subprocess.run(f"docker run ... \{nom\}", shell=True)} laisserait un
utilisateur nommer sa machine \cmd{x; rm -rf /} et faire exécuter n'importe quoi
au serveur. C'est une \emph{injection de commande}.

Passer une \textbf{liste} d'arguments, sans shell, supprime le problème à la
racine : il n'y a aucun interpréteur pour relire la chaîne. Chaque élément de la
liste est un argument, point.
\end{piege}

\apport
La commande imposée est visible et vérifiable. Et les tests peuvent espionner
cette unique fonction pour vérifier ce qui part, sans Docker.

% -----------------------------------------------------------------------------
\outil{Les étiquettes Docker}{brique 23 · l'identité des machines}

\quoi
Des paires clé-valeur attachées à un conteneur à sa création, et sur lesquelles
on peut filtrer.

\pourquoi
C'est la brique née d'un incident. Le Faucheur doit savoir quels conteneurs lui
appartiennent. La première version les reconnaissait à leur nom, qui commence
par \cmd{insacloud\_}.

Puis la pile \cmd{docker compose} est arrivée, avec des conteneurs nommés
\cmd{insacloud-web}, \cmd{insacloud-db}, \cmd{insacloud-faucheur}. Le Faucheur
les a pris pour des machines louées orphelines et les a détruits.
\textbf{Il s'est supprimé lui-même.}

\role
Deux étiquettes, apposées à toute machine louée :

\begin{lstlisting}[language=Python]
MACHINE_LABEL = "insacloud.role"          # = "rented-machine"
CLUSTER_LABEL = "insacloud.cluster"       # = nom du deploiement
\end{lstlisting}

Les deux filtres sont combinés : une machine doit être une machine louée
\textbf{et} appartenir à ce déploiement. Le second a été ajouté après un
deuxième incident : deux déploiements partageant un démon se détruisaient
mutuellement leurs conteneurs.

\begin{retenir}
Un nom est une commodité d'affichage ; une étiquette est une identité. Dès
qu'un programme doit décider de \textbf{détruire} quelque chose, il doit se
fonder sur une marque qu'il a lui-même apposée, jamais sur une convention de
nommage.
\end{retenir}

\apport
Le Faucheur ne peut plus se tromper de cible. Deux plateformes peuvent
cohabiter sur un même démon Docker.

% -----------------------------------------------------------------------------
\outil{Le transport SSH de Docker}{brique 24 · piloter les nœuds distants}

\quoi
Le client Docker sait parler à un démon distant à travers une connexion SSH :
\cmd{docker -H ssh://utilisateur@machine ps}.

\pourquoi
L'autre façon de commander un démon distant est d'exposer son API sur le réseau,
port 2375 en clair ou 2376 en TLS. C'est plus direct, et c'est une très mauvaise
idée : \textbf{qui parle au démon Docker d'une machine en est root}. Une erreur
de pare-feu, et toute la plateforme est compromise.

Le transport SSH ne demande aucun port supplémentaire, réutilise une
authentification par clé déjà éprouvée, et chiffre tout.

\role
Une fonction construit le préfixe de commande selon le nœud visé :

\begin{lstlisting}[language=Python]
def docker_command(worker: str = "local") -> list:
    host = worker_host(worker)
    if host is None:
        return ["docker"]
    return ["docker", "-H", f"ssh://{DOCKER_SSH_USER}@{host}"]
\end{lstlisting}

Conséquence heureuse : \textbf{le code est le même} en mono-hôte et en cluster.
Seul le préfixe change. C'est ce qui permet de développer sur un portable sans
monter quatre machines virtuelles.

\apport
Un parc distribué sans ouvrir un seul port Docker sur le réseau.

% -----------------------------------------------------------------------------
\outil{L'agent des nœuds}{brique 25 · l'inscription spontanée}

\quoi
Un programme Python de 263 lignes qui tourne sur chaque worker. Au démarrage il
se présente au contrôleur ; ensuite il signale sa présence toutes les trente
secondes.

\pourquoi
Sans lui, il faudrait écrire la liste des workers dans la configuration du
contrôleur. Ajouter un nœud deviendrait : modifier un fichier, redéployer. Et un
nœud en panne resterait dans la liste jusqu'à ce qu'un humain l'en retire.

L'agent inverse la responsabilité : c'est le nœud qui se déclare. Le contrôleur
ne contient \textbf{aucune adresse de worker}.

Il n'utilise que la bibliothèque standard --- \cmd{urllib} pour les appels,
\cmd{/proc/meminfo} pour la mémoire --- précisément pour qu'il n'y ait rien à
installer sur les workers.

\role
Deux routes de l'API : \cmd{POST /workers/register} et
\cmd{POST /workers/heartbeat}, protégées par un jeton partagé comparé en temps
constant. L'inscription est retentée avec un recul exponentiel : au démarrage
d'un cluster, un worker peut être prêt avant le contrôleur.

\begin{center}
\small
\begin{tabularx}{\textwidth}{@{}lX@{}}
\toprule
\textbf{État} & \textbf{Signification} \\
\midrule
\cmd{AVAILABLE} & le nœud bat, Docker répond, il reste de la place \\
\cmd{BUSY} & le nœud bat, mais il est plein : on ne lui envoie plus rien \\
\cmd{OFFLINE} & plus de battement depuis 90 s : ses machines sont recréées ailleurs \\
\bottomrule
\end{tabularx}
\end{center}

\begin{retenir}
La distinction \cmd{BUSY} / \cmd{OFFLINE} est essentielle. Un nœud saturé est en
bonne santé : migrer ses machines serait coûteux et inutile. Confondre « plein »
et « mort » déclencherait des reprises à répétition.
\end{retenir}

\apport
Pour ajouter un worker, on l'allume. C'est tout. C'est la Fonctionnalité 7 du
sujet, et c'est ce qui rend le parc élastique.

% -----------------------------------------------------------------------------
\outil{L'ordonnanceur}{brique 26 · choisir le nœud}

\quoi
La fonction qui décide sur quel worker placer une nouvelle machine.

\pourquoi
Les stratégies possibles sont nombreuses : tourniquet, aléatoire, par
pondération de ressources. On a choisi la plus simple qui soit défendable ---
le nœud joignable le \textbf{moins chargé} --- pour une raison précise : elle
est \textbf{déterministe}, donc testable.

\begin{retenir}
Un ordonnanceur dont le départage en cas d'égalité est aléatoire donne un test
intermittent. Et un test intermittent finit toujours par être ignoré. Ici,
l'égalité se tranche par ordre alphabétique.
\end{retenir}

\role
Trois règles : ne considérer que les nœuds \cmd{AVAILABLE} ; parmi eux, prendre
celui qui héberge le moins de machines ; départager par le nom. Puis réserver un
port libre, dont l'unicité est garantie par l'index partiel de la couche 2.

\apport
La charge se répartit, et le comportement est reproductible --- donc
vérifiable par un test.

% -----------------------------------------------------------------------------
\outil{Le Faucheur}{brique 27 · le démon de nettoyage}

\quoi
Un processus séparé qui tourne en boucle, toutes les dix secondes, avec trois
missions : détruire ce qui a expiré, réconcilier Docker et la base, et
surveiller le parc.

\pourquoi
Personne ne viendra cliquer sur « supprimer » à trois heures du matin. Il faut
un service qui surveille les horloges.

Pourquoi un \textbf{processus séparé} plutôt qu'un fil d'exécution dans
l'application ? Parce que Gunicorn lance quatre processus : il y aurait quatre
Faucheurs qui se marcheraient dessus. Un service distinct, démarré une fois par
systemd, évite le problème et peut être redémarré sans couper le site.

\role
Trois tâches à chaque cycle.

\textbf{Expirer} : une seule requête SQL, celle que sert l'index
\cmd{(status, end\_time)}.

\textbf{Réconcilier} : la base dit ce qui \emph{devrait} exister, Docker ce qui
existe \emph{vraiment}. Un conteneur orphelin est supprimé ; une instance sans
conteneur passe à « arrêtée ».

\textbf{Surveiller le parc} : déclarer hors ligne les nœuds muets, et faire
recréer leurs machines ailleurs.

\begin{piege}{Un nœud injoignable n'est pas un nœud vide}
C'est le piège le plus dangereux de la réconciliation. Si le nœud ne répond pas,
\cmd{docker ps} renvoie une liste vide. Un code naïf en conclurait que toutes
ses machines ont disparu et les marquerait « arrêtées » --- alors qu'elles
tournent parfaitement.

D'où la règle : un nœud injoignable est \textbf{ignoré}, jamais interprété.
\end{piege}

\apport
Les machines disparaissent toutes seules à l'échéance, les ports se libèrent, et
la base ne dérive plus de la réalité. C'est la Fonctionnalité 12, et la moitié
de la 13.

% -----------------------------------------------------------------------------
\outil{systemd}{brique 28 · faire tourner les services, et les réparer}

\quoi
Le gestionnaire de services de Linux : il démarre les programmes au bon moment,
les redémarre s'ils meurent, et collecte leurs journaux.

\pourquoi
L'alternative serait de lancer les processus à la main dans un \cmd{screen}, ou
d'écrire des scripts d'init. systemd apporte trois choses qu'on ne veut pas
réécrire : le redémarrage automatique, l'ordre de démarrage selon les
dépendances, et les journaux consultables par \cmd{journalctl}.

\role
Deux unités et un \emph{timer}.

L'unité de l'application lance Gunicorn sous un compte sans privilège, avec les
secrets lus dans un fichier séparé :

\begin{lstlisting}
[Service]
User=insacloud
EnvironmentFile=/etc/insacloud/insacloud.env     # secrets, en 0600
ExecStart=/opt/insacloud/.venv/bin/gunicorn --preload -w 4 ...
Restart=always
\end{lstlisting}

\begin{piege}{Un secret dans une unité systemd est public}
Mettre \cmd{Environment=INSACLOUD\_VAULT\_KEY=...} dans le fichier d'unité le
rend lisible par tous : les unités sont en lecture pour tout le monde, et
\cmd{systemctl show} les affiche. La bonne façon est un fichier séparé en
\cmd{0600}, référencé par \cmd{EnvironmentFile=}.
\end{piege}

\textbf{Le watchdog} est un \emph{timer} qui se déclenche chaque minute. Il
vérifie que nginx répond, que le site renvoie bien du HTTPS, et que le Faucheur
tourne. Il redémarre ce qui est tombé et journalise les workers injoignables.

Pourquoi un timer systemd plutôt que cron ? Parce qu'il s'intègre aux journaux,
connaît l'état du service qu'il surveille, et ne démarre pas avant que ses
dépendances soient prêtes. Cron ne sait rien de tout cela.

\apport
Le service se relance seul après une panne ou un redémarrage du serveur. C'est
le troisième ajout au-delà du sujet.

% -----------------------------------------------------------------------------
\section*{Ce que la plateforme sait faire après la couche 4}

Elle loue, place, surveille et détruit. Elle survit à la panne d'un nœud. Elle
se répare elle-même. Mais tout cela n'existe que sur des machines que
quelqu'un a installées à la main --- et ça, ce n'est pas tenable.

\chapter{Couche 5 — L'infrastructure}

Jusqu'ici, tout suppose des serveurs déjà installés. Cinq briques pour que ces
serveurs se créent et se configurent tout seuls, à l'identique, autant de fois
qu'on veut.

\begin{figure}[htbp]
\centering
\resizebox{0.92\textwidth}{!}{%
\begin{tikzpicture}[font=\small,
  b/.style={draw=grisligne,fill=grisfond,rounded corners=3pt,minimum height=11mm,
            minimum width=3.2cm,align=center},
  f/.style={-{Stealth[length=2.5mm]},gristexte,thick}]
\node[b] (r) at (0,0) {\textbf{Rien}\\\textcolor{gristexte}{\scriptsize un dossier}};
\node[b] (v) at (4.2,0) {\textbf{Vagrant}\\\textcolor{gristexte}{\scriptsize 4 VM nues}};
\node[b] (a) at (8.4,0) {\textbf{Ansible}\\\textcolor{gristexte}{\scriptsize 4 VM configurées}};
\node[b] (p) at (12.6,0) {\textbf{La plateforme}\\\textcolor{gristexte}{\scriptsize qui tourne}};
\draw[f] (r) -- node[above,font=\scriptsize]{\cmd{vagrant up}} (v);
\draw[f] (v) -- node[above,font=\scriptsize]{\cmd{ansible-playbook}} (a);
\draw[f] (a) -- node[above,font=\scriptsize]{démarrage} (p);
\end{tikzpicture}}
\caption{Deux commandes séparent un dossier vide d'une plateforme en
fonctionnement. Tout ce qui se passe entre les deux est décrit dans des fichiers
versionnés.}
\end{figure}

% -----------------------------------------------------------------------------
\outil{Vagrant}{brique 29 · créer les machines virtuelles}

\quoi
Un outil qui décrit des machines virtuelles dans un fichier texte et les crée en
une commande. \cmd{vagrant up} les monte, \cmd{vagrant destroy} les efface.

\pourquoi
L'alternative est de créer les VM à la main dans l'interface de VirtualBox ou de
VMware : choisir l'image, la mémoire, le réseau, installer le système. Une demi-
heure par machine, et personne ne saura ensuite ce qui a été coché.

Terraform aurait pu faire le travail, mais il vise les fournisseurs de cloud
publics ; pour des VM locales, Vagrant est plus direct et c'est l'outil que le
sujet nomme.

\role
Le \fichier{Vagrantfile} décrit un contrôleur en 192.168.56.10 et trois workers
en .11, .12, .13, avec leur mémoire et leurs cœurs. Le nombre de workers et les
tailles sont réglables sans toucher au fichier :

\begin{lstlisting}[language=bash]
WORKERS=2 CTRL_MEM=1536 WORKER_MEM=1280 vagrant up
\end{lstlisting}

\begin{retenir}
Rendre les tailles paramétrables coûte trois lignes et change tout pour qui
reprend le projet : on peut le faire tourner sur un portable modeste. Un fichier
qui n'offre qu'une seule taille finit toujours par être modifié à la main, donc
désynchronisé du dépôt.
\end{retenir}

\apport
Le parc entier se crée en une commande, et se détruit aussi vite. On peut tout
casser pour vérifier que tout se remonte.

% -----------------------------------------------------------------------------
\outil{Les providers : VMware, VirtualBox, Parallels, QEMU}{brique 30 · exécuter les VM}

\quoi
Le \emph{provider} est le logiciel de virtualisation que Vagrant pilote. Vagrant
ne virtualise rien lui-même : il donne des ordres à l'un d'eux.

\pourquoi
Un seul provider aurait suffi... sur une seule architecture. Le projet devait
tourner sur des postes Intel et sur des Mac à puce Apple.

\textbf{VirtualBox} est le provider par défaut et le plus répandu, mais son
support ARM reste expérimental : sur un Mac récent, il ne convient pas.
\textbf{VMware Fusion} fonctionne nativement sur Apple Silicon.
\textbf{Parallels} aussi, mais il est payant. \textbf{QEMU} est libre et
fonctionne partout, au prix de limitations.

Déclarer les quatre laisse le même fichier fonctionner sur n'importe quel poste.

\role
Chaque provider a son bloc dans le \fichier{Vagrantfile}, avec ses réglages
propres.

\begin{piege}{Chaque provider a ses limites}
\cmd{vagrant-qemu} ne sait pas créer de réseau privé : l'accès SSH doit être
redirigé sur \cmd{localhost}. Le fichier le détecte et ajuste la configuration.
Sans cela, le déploiement réussirait et rien ne serait joignable --- la pire
sorte de panne, celle qui ne dit pas son nom.
\end{piege}

\apport
Le projet se monte sur le poste de n'importe quel membre du groupe, et sur
l'ordinateur de la salle de soutenance.

% -----------------------------------------------------------------------------
\outil{Ansible}{brique 31 · configurer les machines}

\quoi
Un outil qui se connecte en SSH aux machines et applique un \textbf{état décrit
en YAML} : tel paquet installé, tel fichier présent avec tels droits, tel
service démarré.

\pourquoi
L'alternative est un script shell. La différence n'est pas cosmétique.

Un script décrit des \textbf{actions} : installe, crée, copie, redémarre. Le
relancer refait tout, écrase, redémarre inutilement, et échoue si quelque chose
existe déjà.

Ansible décrit un \textbf{état voulu}. Il regarde, compare, et n'agit que si
nécessaire. C'est la propriété d'\textbf{idempotence} : appliquer deux fois
donne le même résultat qu'appliquer une fois.

Les concurrents --- Puppet, Chef, Salt --- exigent un agent installé sur chaque
machine. Ansible n'a besoin que de SSH et de Python, déjà présents. C'est aussi
l'outil que le sujet nomme.

\begin{notion}{Idempotent ne veut pas dire « ne fait rien »}
Cela veut dire « vérifie l'état, et n'agit que s'il diffère de celui qui est
voulu ». La nuance compte : c'est elle qui permet de relancer le déploiement
sans crainte, à n'importe quel moment, y compris sur un parc à moitié
configuré.

La preuve se lit dans le compte rendu : au second passage, \cmd{changed=0} sur
tous les nœuds.
\end{notion}

\role
Neuf rôles, quatre playbooks. Chaque rôle est autonome, avec ses variables par
défaut, ses fichiers et ses gabarits.

\begin{center}
\small
\begin{tabularx}{\textwidth}{@{}lX@{}}
\toprule
\textbf{Rôle} & \textbf{Ce qu'il garantit} \\
\midrule
\cmd{security\_hardening} & le pare-feu UFW, avec les ports du groupe \\
\cmd{controller\_key} & le compte de service et sa clé SSH de pilotage \\
\cmd{docker} & Docker Engine sur les workers, le client seul sur le contrôleur \\
\cmd{tls\_cert} & un certificat propre à chaque nœud \\
\cmd{worker} & l'accès SSH restreint du contrôleur \\
\cmd{docker\_images} & les images des machines louées, construites sur place \\
\cmd{worker\_agent} & l'agent d'inscription, en service systemd \\
\cmd{reverse\_proxy} & nginx en HTTPS devant Gunicorn \\
\cmd{webapp} & l'application, le Faucheur, le watchdog, les secrets \\
\bottomrule
\end{tabularx}
\end{center}

\begin{piege}{Un rôle appelé avec des paramètres différents s'exécute deux fois}
Ansible dédoublonne les rôles... sauf si on les appelle avec des paramètres
différents : il les considère alors comme distincts. Le rôle \cmd{docker} était
appelé directement \emph{et} tiré en dépendance par \cmd{docker\_images}, avec
un paramètre différent --- il s'exécutait deux fois, installant puis
désinstallant le démon.

La correction : passer par une \textbf{variable de play}, qui s'applique aux
deux appels et les fait fusionner. Nous l'avons vérifié par l'expérience, en
instrumentant le rôle pour qu'il journalise chaque exécution.
\end{piege}

\apport
Quatre machines nues deviennent une plateforme en fonctionnement, en une
commande, à l'identique. Et le déploiement peut être rejoué sans risque.

% -----------------------------------------------------------------------------
\outil{Ansible Vault}{brique 32 · les secrets dans Git}

\quoi
Un fichier chiffré en AES-256 qu'Ansible déchiffre au moment de jouer le
playbook, avec une phrase de passe fournie à l'exécution.

\pourquoi
Les secrets --- clé de session, clé du coffre, mot de passe de la base, jeton des
agents --- doivent arriver sur le serveur. Trois mauvaises façons de s'y
prendre :

\textbf{Les écrire en clair dans le dépôt} : \cmd{gitleaks} les trouverait, et
n'importe qui ayant accès au dépôt aussi.

\textbf{Les passer en variables d'environnement à la main} : il faut s'en
souvenir, et ils finissent dans l'historique du shell.

\textbf{Les mettre dans un gestionnaire externe} --- HashiCorp Vault, AWS Secrets
Manager : la bonne solution en entreprise, mais c'est un service de plus à
déployer et à maintenir.

Ansible Vault est le compromis juste pour ce projet : un fichier,
\textbf{versionnable tel quel}, illisible sans la phrase de passe.

\role
\fichier{3\_ansible/group\_vars/insacloud/vault.yml} contient trois secrets. Le
rôle \cmd{webapp} vérifie même leur forme avant de déployer --- la clé Fernet
doit faire exactement 44 caractères --- plutôt que de laisser l'application
échouer au démarrage.

La phrase de passe vit dans un fichier \cmd{.vault\_pass} en \cmd{0600}, exclu
du dépôt par le \fichier{.gitignore}.

\apport
Les secrets voyagent avec le code sans jamais être lisibles. Et le déploiement
reste automatisable : pas d'étape manuelle de copie de clés.

% -----------------------------------------------------------------------------
\outil{L'inventaire dynamique}{brique 33 · découvrir les machines}

\quoi
Un script exécutable qui, au lieu de lire une liste figée de machines, interroge
Vagrant et produit l'inventaire en JSON.

\pourquoi
L'inventaire statique est un fichier qui liste les machines et leurs adresses.
Il a deux défauts : il faut le tenir à jour quand le parc change, et il code en
dur le chemin de la clé SSH --- qui dépend du provider.

C'est un piège réel : l'inventaire du projet contenait
\cmd{vagrant\_provider=virtualbox}. En passant à VMware, Ansible cherchait la
clé au mauvais endroit.

\role
\fichier{3\_ansible/inventaire\_vagrant.py} interroge \cmd{vagrant ssh-config} et
reprend le chemin de clé que Vagrant \emph{déclare lui-même}. Il ne découvre que
les machines réellement démarrées.

\begin{lstlisting}[language=bash]
ansible-playbook -i inventaire_vagrant.py site.yml --ask-vault-pass
\end{lstlisting}

\apport
Le même dépôt se déploie sur les quatre providers sans modification, et sur un
parc de taille variable. C'est le quatrième ajout au-delà du sujet.

% -----------------------------------------------------------------------------
\section*{Ce que la plateforme sait faire après la couche 5}

Elle se construit entièrement toute seule : quatre machines virtuelles, leur
système, leurs paquets, leurs certificats, leurs services, ses secrets. Deux
commandes depuis un dossier vide.

Il reste à garantir qu'elle ne se dégrade pas.

\chapter{Couche 6 — La qualité}

La plateforme marche. Reste à garantir qu'elle marchera encore après la
prochaine modification --- et qu'elle ne livre pas de faille connue. Treize
briques, dont onze tournent à chaque envoi de code.

% -----------------------------------------------------------------------------
\outil{pytest}{brique 34 · les tests}

\quoi
Le cadre de test de Python. Une fonction dont le nom commence par \cmd{test\_}
est un test ; une assertion fausse le fait échouer.

\pourquoi
L'alternative dans la bibliothèque standard est \cmd{unittest}, qui impose
d'écrire des classes et des méthodes à noms imposés. pytest accepte des
fonctions simples, et surtout il apporte les \emph{fixtures} : des morceaux de
préparation réutilisables, dont l'un est central ici.

\role
168 tests, en vingt secondes, répartis en six fichiers. Deux mécanismes font
tout le travail.

\textbf{La base vierge à chaque test.} Une fixture marquée \cmd{autouse=True}
efface et recrée la base avant chaque test, sans qu'on ait à y penser.

\begin{piege}{Des tests qui dépendent de leur ordre}
Si un test laisse des données derrière lui, le suivant peut réussir pour la
mauvaise raison --- ou échouer quand on le lance seul. Le symptôme : « ça passe
en entier mais pas tout seul », ou l'inverse. C'est le genre de bug qui fait
abandonner une suite de tests.
\end{piege}

\textbf{L'espion sur Docker.} Les tests ne lancent aucun conteneur :
\cmd{run\_docker} est remplacé par une fonction qui note ce qu'on lui demande.
On vérifie ainsi que la commande imposée par le sujet part bien telle quelle,
sur une machine qui n'a même pas Docker.

\apport
Les tests tournent partout, en vingt secondes. C'est ce qui permet de les lancer
à chaque envoi plutôt qu'une fois par mois.

% -----------------------------------------------------------------------------
\outil{ruff}{brique 35 · l'analyse de style et d'erreurs}

\quoi
Un analyseur qui lit le code sans l'exécuter et signale ce qui est douteux :
import jamais utilisé, variable inconnue, tournure obsolète, motif dangereux.

\pourquoi
Les outils historiques sont \cmd{flake8} pour le style et \cmd{pylint} pour
l'analyse, plus \cmd{isort} pour les imports. Ruff fait les trois, est écrit en
Rust, et tourne en moins d'une seconde sur tout le projet --- contre plusieurs
dizaines pour pylint. Cette vitesse change l'usage : on le lance avant chaque
commit au lieu de l'éviter.

\role
Configuré dans \fichier{pyproject.toml}, avec six familles de règles
sélectionnées. Une seule est désactivée --- la longueur de ligne --- parce que
les commentaires pédagogiques du projet dépassent parfois.

Il vérifie aussi \cmd{target-version = "py310"} : il refuse les tournures que le
Python de la machine cible ne comprendrait pas.

\apport
Le code reste lisible et cohérent sans discussion de style, et les fautes
d'inattention sont attrapées avant l'exécution.

% -----------------------------------------------------------------------------
\outil{bandit}{brique 36 · l'analyse de sécurité du code (SAST)}

\quoi
\emph{Static Application Security Testing} : il cherche dans le code source des
motifs connus pour être dangereux.

\pourquoi
Un relecteur humain oublie. bandit connaît une liste de pièges classiques et ne
se fatigue pas : concaténation SQL, \cmd{shell=True}, secret en dur, algorithme
de hachage obsolète, fichier temporaire prévisible.

\role
Lancé avec \cmd{--severity-level low --confidence-level low} : on veut tout
voir. Le projet ne désactive \textbf{aucune règle globalement} ; les rares cas
légitimes sont annotés ligne à ligne avec leur justification --- par exemple
\cmd{\# nosec B404} sur l'import de \cmd{subprocess}, puisque le sujet impose de
piloter Docker par sa ligne de commande.

\begin{piege}{Le témoin de vie dans \texttt{/tmp}}
Le Faucheur horodate un fichier à chaque cycle, pour que la sonde du conteneur
sache que sa boucle tourne. Nous l'avions mis dans \cmd{/tmp} ; bandit a
signalé la règle B108, et il avait raison.

Le Faucheur tourne aussi en service systemd sur le contrôleur, où \cmd{/tmp} est
partagé avec tous les comptes locaux : un lien symbolique y suffirait à nous
faire horodater un fichier qui ne nous appartient pas. Le témoin vit désormais
dans le dossier de l'application.
\end{piege}

\apport
Une famille entière d'erreurs est attrapée automatiquement, y compris celles
qu'on vient d'introduire sans y penser.

% -----------------------------------------------------------------------------
\outil{pip-audit}{brique 37 · l'audit des dépendances (SCA)}

\quoi
\emph{Software Composition Analysis} : il confronte la liste des dépendances
épinglées aux bases de vulnérabilités publiques.

\pourquoi
On écrit soi-même quelques milliers de lignes, mais on en importe des
centaines de milliers. Une faille dans Flask ou dans Werkzeug est une faille du
projet, et on ne l'apprend pas en relisant son propre code.

L'alternative est \cmd{safety}, dont la base complète est payante. pip-audit est
publié par l'autorité de packaging Python et utilise les bases publiques.

\role
Lancé sur \fichier{requirements.txt} avec \cmd{--strict}. C'est la brique qui
rend l'épinglage des versions (brique 3) exploitable : on ne peut auditer que ce
qu'on connaît exactement.

\begin{retenir}
Ce travail trouve des failles \textbf{sans qu'on ait rien changé}. Werkzeug
3.1.8 est devenu vulnérable (CVE-2026-102598) entre deux exécutions de la
chaîne : notre code n'avait pas bougé, le monde si. Sans ce contrôle
automatique, nous l'aurions appris trop tard.
\end{retenir}

\apport
Les failles des dépendances sont vues le jour de leur publication, pas le jour
de l'incident.

% -----------------------------------------------------------------------------
\outil{gitleaks}{brique 38 · la recherche de secrets}

\quoi
Un outil qui parcourt le dépôt \textbf{et tout son historique} à la recherche de
clés, jetons et mots de passe.

\pourquoi
C'est la brique qui corrige une illusion : effacer un secret dans un commit
ultérieur ne l'efface pas du dépôt. Il reste dans l'historique, consultable par
quiconque clone.

\role
Lancé sur l'historique complet à chaque envoi. C'est ce qui permet au Vault
d'être versionné sans crainte : s'il était un jour commis en clair, la chaîne le
verrait immédiatement.

\apport
Une erreur humaine --- commettre un fichier de configuration contenant un mot de
passe --- est détectée avant qu'elle ne devienne publique.

% -----------------------------------------------------------------------------
\outil{nuclei et dast.sh}{brique 39 · l'analyse dynamique (DAST)}

\quoi
\emph{Dynamic Application Security Testing} : on attaque l'application
\textbf{en fonctionnement}, par le réseau, comme le ferait un attaquant.

\pourquoi
bandit lit le code ; il ne peut pas savoir si l'en-tête de sécurité est
réellement envoyé, ni si une page protégée l'est vraiment une fois déployée.
Seul un test contre le service qui tourne le dit.

\cmd{nuclei} est un scanner générique avec des milliers de signatures.
\fichier{tests/dast.sh} est un script maison qui vérifie ce que nuclei ne
connaît pas : les règles propres à ce projet.

\role
La chaîne démarre l'application avec Gunicorn, puis lance les deux. Dix-huit
contrôles : en-têtes de sécurité, pages protégées inaccessibles sans session,
jetons CSRF exigés, attributs des cookies.

\apport
On vérifie le comportement réel, pas l'intention du code.

% -----------------------------------------------------------------------------
\outil{hadolint et ansible-lint}{brique 40 · l'analyse de l'infrastructure}

\quoi
Deux analyseurs spécialisés : \cmd{hadolint} pour les Dockerfiles,
\cmd{ansible-lint} pour les rôles Ansible.

\pourquoi
Le code d'infrastructure est du code, avec ses propres pièges. Un Dockerfile
peut figer un mot de passe dans une couche, installer une version non épinglée,
ou tourner en root sans raison. Un rôle Ansible peut être non idempotent sans
que cela se voie.

\role
\cmd{hadolint} analyse les quatre Dockerfiles. C'est lui qui a trouvé le mot de
passe par défaut de la brique 16.

\cmd{ansible-lint} tourne au profil \cmd{production}, le plus strict. Il vérifie
notamment que les tâches ont un nom, que les modules sont qualifiés
(\cmd{ansible.builtin.*}), et signale les usages non idempotents.

\begin{piege}{\texttt{command} et \texttt{shell} ne sont jamais idempotents}
Ces deux modules exécutent ce qu'on leur donne sans rien vérifier : ils sont
donc \emph{toujours} signalés \cmd{changed}. Il faut leur apprendre à se taire
avec \cmd{creates:}, \cmd{removes:} ou \cmd{changed\_when:}.
\end{piege}

\apport
L'infrastructure est relue avec la même exigence que le code applicatif.

% -----------------------------------------------------------------------------
\outil{Trivy}{brique 41 · le scan des images}

\quoi
Un scanner qui inventorie les paquets d'une image Docker et cherche leurs
vulnérabilités connues.

\pourquoi
pip-audit couvre les dépendances Python. Mais une image contient aussi un
système de base --- Debian, Ubuntu, Alpine --- avec ses centaines de paquets.
Une faille OpenSSL dans l'image de base est une faille livrée.

\role
La chaîne inventorie \emph{toutes} les failles sans bloquer, puis relance Trivy
avec un verrou : aucune faille \cmd{HIGH} ou \cmd{CRITICAL}
\textbf{corrigeable} ne passe.

\begin{retenir}
Le mot « corrigeable » fait tout. L'option \cmd{--ignore-unfixed} écarte les
failles sans correctif disponible : les refuser ne protégerait de rien et
bloquerait tout, ce qui finirait par faire désactiver le contrôle.
\end{retenir}

Ce verrou a bloqué nos propres images : la base \cmd{docker:28-cli} traînait
163 failles corrigeables liées à sa version de Go. Il a fallu passer à
\cmd{docker:29-cli} --- deux failles au lieu de 163 --- et ajouter un
\cmd{apt-get upgrade} pour OpenSSL.

\apport
Aucune image vulnérable ne part au registre. Et la preuve que le contrôle sert :
il nous a refusé les nôtres.

% -----------------------------------------------------------------------------
\outil{Syft}{brique 42 · l'inventaire de ce qu'on livre (SBOM)}

\quoi
\emph{Software Bill of Materials} : la liste exhaustive des composants d'une
image livrée, avec leurs versions.

\pourquoi
Quand une faille majeure est annoncée, la première question est : « sommes-nous
touchés ? ». Sans inventaire, il faut reconstruire l'image et chercher. Avec un
SBOM archivé, on lit un fichier.

C'est aussi une exigence réglementaire croissante pour tout logiciel livré.

\role
Généré à chaque publication et conservé en artefact de la chaîne --- environ
340 Ko décrivant tout ce que contiennent les images.

\apport
Le projet peut répondre, des mois plus tard, à une question sur une faille
publiée après sa livraison.

% -----------------------------------------------------------------------------
\outil{GitHub Actions}{brique 43 · l'orchestration de tout cela}

\quoi
Le service qui exécute automatiquement une suite de travaux à chaque envoi de
code, sur des machines neuves.

\pourquoi
Les onze contrôles précédents ne valent que s'ils tournent. Lancés à la main,
ils sont oubliés le jour où l'on est pressé --- c'est-à-dire le jour où l'on en
a le plus besoin.

Les alternatives --- GitLab CI, Jenkins --- auraient demandé un serveur à
administrer. GitHub Actions est intégré au dépôt, et son exécution sur une
machine neuve à chaque fois garantit qu'on ne dépend pas d'un reste de
configuration locale.

\role
Onze travaux dans \fichier{.github/workflows/ci.yml}, 601 lignes. Et surtout une
\textbf{barrière} : la publication et le déploiement ne s'exécutent que si les
neuf précédents sont au vert, et seulement sur la branche principale.

\begin{retenir}
Sans dépendance déclarée (\cmd{needs:}), GitHub Actions lance tous les travaux
en parallèle --- et publierait l'image pendant que les tests échouent. C'est la
dépendance explicite qui transforme une liste de contrôles en véritable
barrière.
\end{retenir}

\begin{piege}{La chaîne vérifie sa propre syntaxe}
Une erreur de syntaxe shell dans un bloc \cmd{run:} ne se voit qu'au moment où
l'étape s'exécute, parfois après dix minutes. Trois exécutions ont échoué pour
deux guillemets déséquilibrées.

Un contrôle a été ajouté : il extrait chaque bloc \cmd{run:}, neutralise les
expressions de modèle, et passe le tout à \cmd{bash -n} avant que quoi que ce
soit ne démarre.
\end{piege}

\apport
Onze contrôles qui tournent vraiment, à chaque fois, sur une machine neuve.

% -----------------------------------------------------------------------------
\outil{ghcr.io et docker compose}{brique 44 · livrer et faire tourner}

\quoi
\textbf{ghcr.io} est le registre d'images de GitHub : l'endroit où les images
construites sont publiées. \textbf{docker compose} décrit une pile de
conteneurs dans un fichier et la démarre en une commande.

\pourquoi
Pour le registre : il est intégré au dépôt, l'authentification est celle de
GitHub Actions, et les images sont versionnées avec le commit qui les a
produites.

Pour compose : monter quatre machines virtuelles demande du temps et de la
mémoire. Pour développer, pour démontrer, et pour faire tourner les tests
d'intégration dans la chaîne, on veut la pile complète en une commande.

\role
Les images sont étiquetées trois fois : par version, par commit, et
\cmd{latest}. Le fichier \fichier{docker-compose.yml} décrit la base,
l'application, le Faucheur, et les six constructions d'images sous un profil
séparé.

Deux mécanismes y méritent d'être signalés.

\textbf{L'ancre YAML} : un bloc nommé une fois, repris par l'application et le
Faucheur. Il devient impossible qu'ils divergent sur une variable.

\textbf{Le refus de démarrer sans secret} : la syntaxe
\cmd{\$\{POSTGRES\_PASSWORD:?manquant\}} fait échouer le démarrage si la
variable n'est pas définie. L'erreur arrive tout de suite, pas trois semaines
plus tard.

\begin{piege}{Un volume par service là où il en fallait un partagé}
Le Faucheur \emph{écrit} les archives, l'application les \emph{sert} au
téléchargement. Avec un volume par service, l'utilisateur verrait dans son
tableau de bord une archive que le serveur ne trouverait pas. Le même volume
nommé doit être monté dans les deux.
\end{piege}

\gros[0.86]{zoom-archives.png}{Les archives servies par l'application, écrites
par le Faucheur : deux services, un seul volume.}

\apport
Le projet se livre et se lance en une commande, en développement comme en
intégration continue.

% -----------------------------------------------------------------------------
\section*{Ce que la plateforme sait faire après la couche 6}

Elle se vérifie elle-même. Onze contrôles à chaque envoi, dont six de sécurité,
et aucune livraison si l'un échoue. Elle a trouvé dix défauts réels que la
relecture n'avait pas vus.

\chapter{Récapitulatif}

\section{Les quarante-quatre briques, d'un coup d'œil}

\begin{center}
\scriptsize
\begin{tabularx}{\textwidth}{@{}clXX@{}}
\toprule
\textbf{\#} & \textbf{Outil} & \textbf{Son rôle ici} & \textbf{Écarté au profit de lui} \\
\midrule
\multicolumn{4}{@{}l}{\textbf{\color{insared}Couche 0 — Le socle}} \\
1 & Git & historique, branches, base de la CI & copier le dossier \\
2 & Python 3.10 & tout le code, version de la cible & Node, Go (agent à déployer) \\
3 & venv + épinglage & versions identiques partout & installation globale \\
\addlinespace
\multicolumn{4}{@{}l}{\textbf{\color{insared}Couche 1 — L'application web}} \\
4 & Flask & routes HTML et JSON, sessions & Django (trop de boîtes noires) \\
5 & Jinja2 & gabarits, échappement automatique & concaténation de chaînes \\
6 & Werkzeug & hachage pbkdf2, ProxyFix & hachage écrit à la main \\
7 & Gunicorn & 4 ouvriers, production & serveur de dév. de Flask \\
8 & nginx & TLS, statiques, en-têtes & Apache, Caddy \\
9 & Tailwind précompilé & apparence sans CDN & Tailwind par CDN (interdit par la CSP) \\
\addlinespace
\multicolumn{4}{@{}l}{\textbf{\color{insared}Couche 2 — Les données}} \\
10 & SQLAlchemy & 7 tables, pas d'injection SQL & SQL à la main \\
11 & Alembic & migrations versionnées & \cmd{ALTER TABLE} manuels \\
12 & SQLite & développement et tests & un serveur à lancer \\
13 & PostgreSQL & production, index partiels & MySQL \\
14 & Fernet & coffre des mots de passe root & chiffrement maison \\
\addlinespace
\multicolumn{4}{@{}l}{\textbf{\color{insared}Couche 3 — Les machines louées}} \\
15 & Docker & machines en moins d'une seconde & machines virtuelles \\
16 & Dockerfile & 6 images reproductibles & image figée à la main \\
17 & supervisord & 5 services, PID 1 correct & script shell, systemd \\
18 & OpenSSH & l'accès imposé par le sujet & --- \\
19 & ttyd & terminal sans rien installer & Shellinabox, Wetty \\
20 & Xvnc + XFCE + noVNC & bureau dans le navigateur & GNOME, KDE (trop lourds) \\
21 & entrypoint.sh & assemble la machine selon le mode & six images distinctes \\
\addlinespace
\multicolumn{4}{@{}l}{\textbf{\color{insared}Couche 4 — Le pilotage}} \\
22 & subprocess & la commande imposée, visible & SDK docker-py \\
23 & étiquettes Docker & identité des machines louées & préfixe de nom \\
24 & \cmd{docker -H ssh://} & piloter sans ouvrir de port & API Docker exposée \\
25 & l'agent & inscription et battement de cœur & liste de workers en dur \\
26 & l'ordonnanceur & nœud le moins chargé, déterministe & tirage aléatoire \\
27 & le Faucheur & expiration, réconciliation, reprise & fil dans l'application \\
28 & systemd + watchdog & services relancés, auto-réparation & cron, lancement manuel \\
\addlinespace
\multicolumn{4}{@{}l}{\textbf{\color{insared}Couche 5 — L'infrastructure}} \\
29 & Vagrant & 4 VM en une commande & création à la main \\
30 & VMware et 3 autres & fonctionne sur Intel et Apple & VirtualBox seul \\
31 & Ansible & 9 rôles idempotents & scripts shell \\
32 & Ansible Vault & secrets versionnables & secrets en clair, Vault externe \\
33 & inventaire dynamique & découvre les VM et leur clé & inventaire figé \\
\addlinespace
\multicolumn{4}{@{}l}{\textbf{\color{insared}Couche 6 — La qualité}} \\
34 & pytest & 168 tests en 20 s, sans Docker & unittest \\
35 & ruff & style et erreurs, en moins d'1 s & flake8 + pylint + isort \\
36 & bandit & motifs dangereux (SAST) & relecture humaine \\
37 & pip-audit & CVE des dépendances (SCA) & safety (base payante) \\
38 & gitleaks & secrets, historique compris & recherche manuelle \\
39 & nuclei + dast.sh & l'application en marche (DAST) & analyse statique seule \\
40 & hadolint + ansible-lint & l'infrastructure relue & relecture humaine \\
41 & Trivy & CVE des images, verrou & publication sans scan \\
42 & Syft & SBOM conservé & aucun inventaire \\
43 & GitHub Actions & 11 travaux, barrière & GitLab CI, Jenkins \\
44 & ghcr.io + compose & livrer, et lancer en une commande & registre tiers \\
\bottomrule
\end{tabularx}
\end{center}

\section{Qui réalise quoi dans le cahier des charges}

\begin{center}
\scriptsize
\begin{tabularx}{\textwidth}{@{}clX@{}}
\toprule
\textbf{\#} & \textbf{Fonctionnalité} & \textbf{Briques mobilisées} \\
\midrule
1 & Préparation, environnement & Git · Python · venv \\
2 & API Web Flask & Flask · Jinja2 · Werkzeug · Gunicorn · nginx · Tailwind \\
3 & Base de données & SQLAlchemy · Alembic · SQLite · PostgreSQL \\
4 & Conteneurisation Docker & Docker · Dockerfile · supervisord · entrypoint \\
5 & Infrastructure Vagrant & Vagrant · les quatre providers \\
6 & Automatisation Ansible & Ansible · Vault · inventaire dynamique \\
7 & Worker Agent & l'agent · les étiquettes · l'API des nœuds \\
8 & Resource Manager & l'ordonnanceur · l'index partiel de PostgreSQL \\
9 & Service de location & Flask · services.py · SQLAlchemy \\
10 & Déploiement des instances & subprocess · Docker · \cmd{docker -H ssh://} \\
11 & Accès SSH & OpenSSH · ttyd · Xvnc + XFCE + noVNC \\
12 & Gestion de l'expiration & le Faucheur · l'index \cmd{(status, end\_time)} \\
13 & Haute disponibilité & supervisord · \cmd{restart=always} · live-restore · l'agent · le Faucheur \\
14 & Sécurité DevSecOps & Fernet · bandit · pip-audit · gitleaks · nuclei · Trivy · Vault \\
15 & CI/CD & GitHub Actions · ghcr.io · Syft \\
16 & Tests d'intégration & pytest · docker compose · integration.sh \\
17 & Documentation & les schémas engendrés · ce document · le guide de 93 pages \\
\bottomrule
\end{tabularx}
\end{center}

\begin{retenir}
Aucune fonctionnalité n'est portée par un seul outil. La haute disponibilité,
par exemple, naît de cinq briques posées dans trois couches différentes :
supervisord dans la machine, \cmd{--restart=always} et \cmd{live-restore} sur le
nœud, l'agent et le Faucheur dans le parc. C'est pour cela qu'on ne peut pas
« ajouter la résilience à la fin ».
\end{retenir}

\section{Les sept ajouts au-delà du sujet}

\begin{center}
\small
\begin{tabularx}{\textwidth}{@{}XX@{}}
\toprule
\textbf{Ajout} & \textbf{Briques qui le rendent possible} \\
\midrule
Parc distribué à trois workers, alors qu'un hôte suffisait &
l'agent · l'ordonnanceur · \cmd{docker -H ssh://} \\
\addlinespace
Terminal web, en plus du SSH demandé &
ttyd · supervisord \\
\addlinespace
Bureau graphique complet dans le navigateur &
Xvnc · XFCE · noVNC · supervisord \\
\addlinespace
Coffre à secrets réaffichable après re-authentification &
Fernet · les sessions Flask · Ansible Vault \\
\addlinespace
Watchdog d'auto-réparation du service &
systemd · un timer d'une minute \\
\addlinespace
Inventaire dynamique indépendant du provider &
Vagrant · un script Python de 140 lignes \\
\addlinespace
Sauvegarde du travail avant destruction &
Docker \cmd{exec} · \cmd{tar} en flux · une table de plus \\
\bottomrule
\end{tabularx}
\end{center}

\section{Ce qu'il faut retenir}

\begin{description}
\item[Un outil se comprend par ce qu'on a écarté.] Savoir qu'on utilise Gunicorn
  ne vaut rien ; savoir que le serveur de Flask traite une requête à la fois,
  si. Chaque fiche de ce document a une section « pourquoi celui-là » pour cette
  raison.

\item[Les couches ne sont pas interchangeables.] On ne peut pas conteneuriser
  avant d'avoir une application, ni automatiser le déploiement avant d'avoir
  quelque chose à déployer. L'ordre dans lequel les briques ont été posées n'est
  pas un récit : c'est une contrainte.

\item[Les outils de la couche 6 ne servent à rien s'ils ne bloquent pas.]
  Un contrôle qui affiche un avertissement que personne ne lit est un coût sans
  bénéfice. Ici, aucune image n'atteint le registre si l'un des onze travaux
  échoue --- et c'est arrivé à nos propres images.

\item[Les meilleures briques viennent d'une panne.] Les étiquettes Docker
  existent parce que le Faucheur s'est supprimé lui-même. Le témoin de vie
  existe parce que la sonde du conteneur ne pouvait pas passer. L'inventaire
  dynamique existe parce que l'inventaire figé codait le provider en dur. Un
  projet sans incident n'a pas de garde-fous.
\end{description}

\vfill
\begin{center}
\color{gristexte}\small
Dépôt : \url{https://github.com/aziz3r/insacloud}\\
Pour la construction du projet pas à pas, voir le guide technique de 93 pages.
\end{center}


\end{document}
